\section{Discussion} \label{sec:discussion}

\todoinline{What else? Maybe discuss the option to reduce num of channels while preserving accuracy or convergence rate?}

\subsection{Why Whitening is Applied Across Features}
\todoinline{Is this trivial info? Consider removing}

Batch Whitening is applied in the \emph{feature} (channel) dimension rather than across tokens (in LLMs) or pixels (in CNNs).  
In both domains, the model relies on correlations \emph{along the spatial or temporal axis}:  
CNNs depend on local pixel relationships, and transformers depend on relationships between tokens.  
Whitening across these axes would explicitly remove the very correlations the architecture is designed to model, thereby degrading attention patterns in LLMs and disrupting spatial structures in vision models.

In contrast, whitening across \emph{feature channels} operates within each token or pixel representation.  
This preserves structural relationships (token–token or pixel–pixel), while decorrelating hidden channels to improve conditioning, stabilize optimization, and reduce feature redundancy.  
Thus, feature-domain whitening complements both convolutional and attention-based models: it improves the representation quality without disturbing the core correlation patterns these architectures are meant to capture.